% ===== PRÉAMBULE CANONIQUE — à coller VERBATIM en tête de CHAQUE .tex =====
\documentclass[11pt,a4paper]{article}
\usepackage[utf8]{inputenc}\usepackage[T1]{fontenc}\usepackage{lmodern}
\usepackage[french]{babel}
\usepackage{amsmath,amssymb,mathtools}\usepackage{icomma}
\usepackage{siunitx}\sisetup{locale=FR}  % \qty{}{}, \si{}, \ang{}
\usepackage{xcolor}\usepackage{colortbl}
\usepackage{tikz}\usepackage{pgfplots}\pgfplotsset{compat=1.18}
\usepackage[siunitx,europeanresistors]{circuitikz} % circuits (élec.)
\usepackage[most]{tcolorbox}\usepackage{enumitem}\usepackage{array,booktabs}
\usepackage[a4paper,margin=2.2cm]{geometry}
\usetikzlibrary{arrows.meta,calc,angles,quotes,positioning,patterns,%
  backgrounds,shapes.geometric,decorations.pathmorphing,decorations.markings,intersections}
\definecolor{primary}{RGB}{222,49,99}\definecolor{primarydark}{RGB}{150,22,55}
\definecolor{defcol}{RGB}{0,114,178}\definecolor{methcol}{RGB}{52,92,148}
\definecolor{excol}{RGB}{0,135,90}\definecolor{attncol}{RGB}{200,40,55}
\tcbset{boxbase/.style={enhanced,breakable,boxrule=0.9pt,arc=2.5pt,
  left=3mm,right=3mm,top=2.3mm,bottom=2.3mm,fonttitle=\bfseries,coltitle=white}}
% --- boîtes TUTORIEL (noms EXACTS, ne pas renommer) ---
\newtcolorbox{cle}[1][]{boxbase,colback=defcol!6,colframe=defcol,
  title={\textsc{À retenir}\ifx\relax#1\relax\relax\else\textmd{ : \itshape #1}\fi}}
\newtcolorbox{methode}[1][]{boxbase,colback=methcol!7,colframe=methcol,sharp corners,
  title={\textsc{Méthode}\ifx\relax#1\relax\relax\else\textmd{ --- \itshape #1}\fi}}
\newtcolorbox{exemplebox}[1][]{boxbase,colback=excol!5,colframe=excol,
  title={\textsc{Exemple}\ifx\relax#1\relax\relax\else\textmd{ : \itshape #1}\fi}}
\newtcolorbox{piege}[1][]{boxbase,colback=attncol!6,colframe=attncol,
  title={\textsc{Piège}\ifx\relax#1\relax\relax\else\textmd{ : \itshape #1}\fi}}
\newtcolorbox{remarque}[1][]{boxbase,colback=black!4,colframe=black!45,
  title={\textsc{Remarque}\ifx\relax#1\relax\relax\else\textmd{ : \itshape #1}\fi}}
% --- boîtes EXERCICE / CORRIGÉ (noms EXACTS, auto-comptées) ---
\newtcolorbox[auto counter]{exo}[1][]{boxbase,colback=primary!4,colframe=primary,
  title={\textsc{Exercice}~\thetcbcounter\ifx\relax#1\relax\relax\else\textmd{ --- \itshape #1}\fi}}
\newtcolorbox[auto counter]{corrige}[1][]{boxbase,colback=excol!4,colframe=excol,
  title={\textsc{Corrigé}~\thetcbcounter\ifx\relax#1\relax\relax\else\textmd{ --- \itshape #1}\fi}}
\newcommand{\vv}[1]{\vec{#1}}\newcommand{\ds}{\displaystyle}
\newcommand{\R}{\mathbb{R}}\newcommand{\N}{\mathbb{N}}\newcommand{\Z}{\mathbb{Z}}
% (un \usepackage{fancyhdr}+en-tête est possible ; il est ignoré côté web)
% =========================================================================
\begin{document}
\begin{corrige}[champ d'un fil rectiligne]
Formule du fil : $B=\dfrac{\mu_0 I}{2\pi d}=\dfrac{2\cdot10^{-7}\,I}{d}$, avec $d=\qty{0.05}{\metre}$ :
\[ B=\frac{2\cdot10^{-7}\times10}{0{,}05}=\frac{2\cdot10^{-6}}{0{,}05}=\qty{4e-5}{\tesla}. \]
(Soit environ le champ magnétique terrestre.)
\end{corrige}

\begin{corrige}[champ au centre d'une spire]
Formule de la spire ($N=1$) : $B=\dfrac{\mu_0 I}{2r}$, avec $r=\qty{0.02}{\metre}$ :
\[ B=\frac{4\pi\times10^{-7}\times5}{2\times0{,}02}
    =\frac{1{,}26\cdot10^{-6}\times5}{0{,}04}
    =\frac{6{,}3\cdot10^{-6}}{0{,}04}\approx\qty{1.6e-4}{\tesla}. \]
\end{corrige}

\begin{corrige}[champ dans un solénoïde]
Formule du solénoïde : $B=\dfrac{\mu_0 N I}{L}$, avec $L=\qty{0.5}{\metre}$ :
\[ B=\frac{4\pi\times10^{-7}\times1000\times3}{0{,}5}
    =\frac{1{,}26\cdot10^{-6}\times3000}{0{,}5}
    =\frac{3{,}77\cdot10^{-3}}{0{,}5}\approx\qty{7.5e-3}{\tesla}=\qty{7.5}{\milli\tesla}. \]
\end{corrige}

\begin{corrige}[reconnaître un électro-aimant]
\begin{enumerate}[label=\alph*)]
  \item C'est un \textbf{électro-aimant} : une bobine parcourue par un courant, autour d'un noyau de
        fer doux, se comporte comme un aimant (une face Nord, une face Sud) tant que le courant circule.
  \item En ouvrant l'interrupteur, le courant s'annule ($I=0$) : le champ du solénoïde disparaît. Le
        fer étant \textbf{doux}, il \textbf{perd immédiatement son aimantation} — l'électro-aimant
        « lâche » sa charge. (Un fer « dur » resterait aimanté : ce serait un aimant permanent.)
\end{enumerate}
\end{corrige}

\begin{corrige}[la décroissance en $1/d$]
Comme $B\propto 1/d$ (à $I$ fixé), le produit $B\times d$ est constant.
\begin{enumerate}[label=\alph*)]
  \item On \textbf{double} la distance $\Rightarrow$ $B$ est \textbf{divisé par 2} :
        $B'=\dfrac{6\cdot10^{-5}}{2}=\qty{3e-5}{\tesla}$.
  \item On \textbf{divise} la distance par 2 $\Rightarrow$ $B$ est \textbf{doublé} :
        $B''=2\times6\cdot10^{-5}=\qty{1.2e-4}{\tesla}$.
\end{enumerate}
\end{corrige}

\end{document}
